\section{Limitations, Ethics, and Reuse}
\label{sec:limitations}

This benchmark should be interpreted with several constraints in mind. First, the main label is a finite-horizon continuation-policy-specific counterfactual target. It is decision-aligned for the benchmarked eviction choice, but it is not an offline-optimal label for the full future control problem. Good performance on this label therefore does not automatically imply online deployment gains.

Second, the labels depend on the continuation policy used after the forced eviction. A different continuation policy or a different horizon can change the numerical target, the optimal set within a decision, and the induced pairwise preferences. Comparisons should therefore be made only across artifacts that keep these label-generation choices fixed.

Third, the open subset is intentionally narrower than the larger internally generated dataset mix. The current family registry still marks some families as pending redistribution review, and two families remain excluded from the open subset. Not every selected family appears in every split in the preserved release. This release-governance boundary is part of the benchmark story and should be made explicit rather than hidden in artifact packaging details.

Fourth, benchmark reuse should preserve provenance distinctions. The release exposes generated supervision and derived tasks, but it does not erase the attribution, citation, or licensing requirements of the upstream raw traces from which candidate rows were derived. Those raw traces remain external artifacts with their own governance constraints.

Fifth, the preserved open release is in a technically strong but not yet fully published state. Lightweight metadata repair, bundle validation, repository tests, and full real-release validation have passed, but anonymous artifact packaging and final public hosting are still pending. The paper should describe this candidly rather than overstating release availability.

Finally, the generated labels are intended for offline benchmarking claims. They can support comparative modeling studies, but they should not be presented as automatic evidence that a learned policy is ready for operational deployment without separate online evaluation, systems integration, and safety review.
